\section{Mukai fixed loci and the Borcherds weight}
\label{sec:mukai-fixed-loci-borcherds-weight}

\providecommand{\cO}{\mathcal{O}}
\providecommand{\K}{\mathrm{K3}}
\providecommand{\Fix}{\mathrm{Fix}}
\providecommand{\Aut}{\mathrm{Aut}}
\providecommand{\symp}{\mathrm{symp}}

\paragraph{Borcherds weight theorem.}
Let \(g_N\in\Aut_{\symp}(\K)\) have order
\[
  N\in\{1,2,3,4,6\},
\]
and let \(\phi_N=\phi^{(g_N)}_{0,1}(\tau,z)\) be the CHL K3 Jacobi
input in Eichler--Zagier normalisation, with
\(q=e^{2\pi i\tau}\), \(y=e^{2\pi iz}\), and
\[
  \phi_1(\tau,z)=y^{-1}+10+y+O(q).
\]
Write \(c_N(0)\) for the \(q^0y^0\)-coefficient of \(\phi_N\), and
write \(\Phi_N\) for the Gritsenko--Nikulin CHL Borcherds product. Then
\[
  \kappa_{\mathrm{BKM}}(\Phi_N)
  =
  \operatorname{wt}(\Phi_N)
  =
  \frac{c_N(0)}{2}.
\]
The primitive CHL constants and weights are
\[
  (c_1(0),c_2(0),c_3(0),c_4(0),c_6(0))
  =
  (10,8,6,4,2),
\]
\[
  (\kappa_{\mathrm{BKM}}(\Phi_1),
   \kappa_{\mathrm{BKM}}(\Phi_2),
   \kappa_{\mathrm{BKM}}(\Phi_3),
   \kappa_{\mathrm{BKM}}(\Phi_4),
   \kappa_{\mathrm{BKM}}(\Phi_6))
  =
  (5,4,3,2,1).
\]
Borcherds 1998 Theorem~13.3 supplies the weight formula.
Gritsenko--Nikulin 1995 Part~II Theorem~2.1 and Gritsenko 1999 compute
the CHL paramodular cusp forms of weights \((5,4,3,2,1)\).

\paragraph{Holomorphic Lefschetz obstruction.}
The ratio identity
\[
  \frac{c_N(0)}{c_1(0)}
  =
  \frac{\chi_{g_N}(\cO_\K)}{\chi(\cO_\K)}
\]
fails for every non-identity CHL order. The cohomology of the structure
sheaf of a K3 surface is
\[
  H^0(\K,\cO_\K)\cong\mathbb C,\qquad
  H^1(\K,\cO_\K)=0,\qquad
  H^2(\K,\cO_\K)\cong H^{0,2}(\K).
\]
A symplectic automorphism fixes the holomorphic two-form, hence fixes
the conjugate class in \(H^{0,2}(\K)\). Therefore
\[
  \chi_{g_N}(\cO_\K)
  =
  \operatorname{tr}\!\left(g_N\mid H^0(\K,\cO_\K)\right)
  +
  \operatorname{tr}\!\left(g_N\mid H^2(\K,\cO_\K)\right)
  =
  2
\]
for every symplectic \(g_N\). Atiyah--Bott 1968 Theorem~4.12 gives the
same value from the local holomorphic Lefschetz residues. Hence
\[
  \frac{\chi_{g_N}(\cO_\K)}{\chi(\cO_\K)}=1,
  \qquad
  \frac{c_N(0)}{c_1(0)}
  =
  \left(1,\frac45,\frac35,\frac25,\frac15\right)
\]
for \(N=(1,2,3,4,6)\).

\paragraph{Mukai fixed loci.}
Mukai 1988 Proposition~1.2 and Nikulin 1980 give the fixed-point counts
for finite symplectic K3 automorphisms. On the CHL orders the
topological Lefschetz number is
\[
  L_{\mathrm{top}}(g_N)
  =
  \sum_i(-1)^i\operatorname{tr}\!\left(g_N\mid H^i(\K,\mathbb C)\right)
  =
  \chi(\K^{g_N})
  =
  (24,8,6,4,2)
\]
for \(N=(1,2,3,4,6)\). The four non-identity fixed-locus counts equal
the four non-identity constants \(c_N(0)\). The identity row separates
the two constructions:
\[
  \chi(\K)=24,\qquad c_1(0)=10.
\]
Thus Mukai fixed loci explain a numerical coincidence on
\(N=2,3,4,6\). The CHL Borcherds-input constant at \(N=1\) comes from
the Jacobi-form normalisation, and the holomorphic Euler trace of
\(\cO_\K\) remains \(2\).

\paragraph{Proof.}
Borcherds 1998 Theorem~13.3 identifies the automorphic weight of a
Borcherds product with one half of the constant term of the
zero-component of the vector-valued modular input:
\[
  \operatorname{wt}(\operatorname{Bor}(\phi_N))=\frac{c_N(0)}2.
\]
By definition, \(\kappa_{\mathrm{BKM}}\) is the denominator-product
weight attached to the corresponding Borcherds--Kac--Moody algebra.
Gritsenko--Nikulin compute the CHL products
\[
  \Phi_N=\Delta^{(N)}_{k(N)},\qquad
  k(N)=(5,4,3,2,1),
\]
for \(N=(1,2,3,4,6)\). Hence \(c_N(0)=2k(N)\), giving
\[
  c_N(0)=(10,8,6,4,2),\qquad
  \kappa_{\mathrm{BKM}}(\Phi_N)=(5,4,3,2,1).
\]
For \(N=1\), \(\Phi_1=\Delta_5\). The monic denominator used in the
Oberdieck--Pandharipande scalar convention is
\[
  D_5=64^{-1}\Delta_5,
\]
and this scalar normalisation leaves
\(\operatorname{wt}(\Delta_5)=\kappa_{\mathrm{BKM}}(\Delta_5)=5\).

\paragraph{Catalogue boundary.}
The primitive CHL ladder and the Gritsenko--Cl\'ery diagonal-divisor
catalogue are distinct automorphic families. The CHL ladder has
\[
  c_N(0)=(10,8,6,4,2)
  \qquad (N=1,2,3,4,6).
\]
The Gritsenko--Cl\'ery catalogue has eight triples
\[
  (t,N;k)=
  (1,1;5),(2,1;2),(1,2;3),(3,1;1),
  (1,3;2),(4,1;\tfrac12),(1,4;\tfrac32),(2,2;1),
\]
with twice-weight tuple
\[
  (10,4,6,2,4,1,3,2).
\]
The common entry is \(F_{1,1}=\Delta_5\). Every other row has its own
level, multiplier system, Jacobi input, and denominator data. The
categorical problem left by the obstruction above is precise:
construct an equivariant Hochschild or Hall-theoretic trace whose
value is \(c_N(0)\). The ordinary holomorphic Euler trace of
\(\cO_\K\) is the constant \(2\).

\paragraph{K3 \texorpdfstring{$\times$}{x} E invariants.}
The CHL ladder contributes the automorphic invariant
\(\kappa_{\mathrm{BKM}}(\Phi_N)=c_N(0)/2\). For \(K3\times E\) the
separate Vol~III invariants are
\[
  \kappa_{\mathrm{cat}}(K3\times E)=0,\qquad
  \kappa_{\mathrm{ch}}(K3\times E)=0,\qquad
  \kappa_{\mathrm{ch}}^{\mathrm{Heis}}(K3\times E)=3,
\]
\[
  \kappa_{\mathrm{BKM}}(\Delta_5)=5,\qquad
  \kappa_{\mathrm{fiber}}(K3\times E)=24.
\]
The value \(24\) is the Mukai-lattice rank of the K3 fibre. The values
\(0,0,3,5,24\) come from distinct constructions.
